\chapter{Puberty}\label{ch:g8:puberty}

Somewhere between about nine and sixteen, every human body
rebuilds itself: the second growth burst of
\cref{rem:g3:stages-of-life:bursts}, and much more than
height. \Cref{rem:g5:human-reproduction:adults} promised
this chapter years ago. Here it is --- the facts, the
timetable's famous unfairness, and the machinery, which will
hand us a discovery bigger than \omterm{def:g8:puberty:puberty}{puberty} itself: the body has
a second communication system.

\section{What changes}

\begin{definition}[Puberty]\label{def:g8:puberty:puberty}
\emph{Puberty}\index{puberty} is the stage in which a
child's body becomes an adult's --- above all, an adult's
in \cref{ch:g8:sexual-reproduction}'s sense: the
reproductive organs mature and begin producing
\emph{\omterm{def:g8:sexual-reproduction:gametes}{gametes}}, and the body takes on its adult form.
It is the biological core of \omterm{def:g3:stages-of-life:stages}{adolescence}
(\cref{def:g3:stages-of-life:stages}).
\end{definition}

\begin{proposition}[The visible changes]\label{prop:g8:puberty:changes}
Alongside the growth burst, over several years:
\begin{itemize}
  \item in girls, commonly from about nine to thirteen:
        breasts develop, hips widen, body hair appears ---
        and the \emph{first menstrual periods} mark the
        cycles of \cref{ch:g8:reproductive-systems} starting
        up;
  \item in boys, commonly a year or two later: the voice
        deepens (``breaks''), \omterm{def:g1:my-body:joint}{shoulders} broaden, body and
        facial hair appear, and the testes begin producing
        \omterm{def:g4:animal-reproduction:sexual}{sperm cells};
  \item in both: \omterm{def:g1:the-five-senses:sense}{skin} turns oilier (the age of spots),
        sweat changes, moods ride the rebuilding --- the
        \omterm{def:g5:brain-in-command:system}{brain}, too, is being rewired toward its adult form.
\end{itemize}
\end{proposition}
\admitted

\begin{remark}[The timetable is personal]\label{rem:g8:puberty:timetable}
\Cref{rem:g1:my-body:different} said it of height;
\omterm{def:g8:puberty:puberty}{puberty} repeats it louder. Classmates of one age can stand
years apart on this road --- early and late are both
normal, the order of changes varies, and nearly everyone
lands in the same place by the end. The ranges above are
wide, and the widest thing about them is how little they
say about any single person.
\end{remark}

\section{The machinery: a second messenger system}

\begin{definition}[Hormones, first encounter]\label{def:g8:puberty:hormones}
\omterm{def:g8:puberty:puberty}{Puberty} is not decided in the organs it changes. It is
\emph{ordered} --- by chemical messengers called
\emph{hormones}\index{hormone}: substances released into
the \emph{\omterm{prop:g4:journey-of-food:absorb}{blood}} by specialized organs, carried
everywhere, and obeyed by the organs built to hear them.
At \omterm{def:g8:puberty:puberty}{puberty}, a control center at the \omterm{def:g5:brain-in-command:system}{brain}'s base wakes,
its hormones instruct the reproductive organs, and
\emph{their} hormones --- carried body-wide by
\cref{prop:g4:heart-and-blood:carries}'s delivery service
--- order the whole rebuild: growth, hair, voice, cycles.
\end{definition}

\begin{example}[One signal, many building sites]\label{ex:g8:puberty:sites}
The \omterm{prop:g4:journey-of-food:absorb}{blood} carries a \omterm{def:g8:puberty:hormones}{hormone} everywhere at once --- which
is why \omterm{def:g8:puberty:puberty}{puberty} is a \emph{coordinated} rebuild: voice,
\omterm{def:g1:my-body:joint}{shoulders} and hair follow the same messenger, on one
schedule, though the sites lie half a body apart. No
\omterm{def:g5:brain-in-command:system}{nerve} wiring connects a voice box to a beard; the
bloodstream's broadcast does
(\cref{ch:g8:hormonal-communication} makes this system a
chapter of its own).
\end{example}

\begin{remark}[Why the mood weather?]\label{rem:g8:puberty:moods}
The same messengers reach the \omterm{def:g5:brain-in-command:system}{brain}, which is itself
mid-renovation: mood swings, new intensities, \omterm{prop:g1:healthy-habits:needs}{sleep}
drifting later --- weather, not fault. Knowing the
machinery helps on both sides of it: the storms are
chemistry plus construction, they are normal, and they
pass.
\end{remark}

\section{Living through the rebuild}

\begin{proposition}[Care during construction]\label{prop:g8:puberty:care}
A body mid-rebuild has a builder's needs, all met before
in this book:
\begin{itemize}
  \item material and fuel: the burst is paid in matter ---
        balanced eating
        (\cref{met:g3:balanced-diet:check}), building
        matter included
        (\cref{exo:g7:digestion-and-nutrients:11});
  \item \omterm{prop:g1:healthy-habits:needs}{sleep}: the \omterm{def:g5:brain-in-command:system}{brain}'s renovation and the body's
        growth both run at night
        (\cref{prop:g5:brain-in-command:protect}) --- and
        the drifted-late clock still needs its hours;
  \item movement: \omterm{def:g3:bones-and-muscles:skeleton}{bones} and \omterm{def:g4:heart-and-blood:heart}{heart} in their last, best
        building years
        (\cref{met:g3:bones-and-muscles:care});
  \item and no sabotage: the substances of
        \cref{prop:g5:healthy-choices:substances} strike
        hardest at bodies and \omterm{def:g5:brain-in-command:system}{brains} still under
        construction.
\end{itemize}
\end{proposition}
\admitted

\begin{example}[Questions have addresses]\label{ex:g8:puberty:questions}
\omterm{def:g8:puberty:puberty}{Puberty} raises questions --- about timing, about bodies,
about feelings --- and they have proper addresses: school
nurses and doctors answer them daily, confidentially and
without surprise; trusted adults have all been through
it; and this book's chapters are built to be reread.
What deserves no trust: playground legend and the
internet's worst corners, which trade on exactly this
age's uncertainties.
\end{example}

\section{Exercises}

\begin{exercise}[$\star$]\label{exo:g8:puberty:1}
Define \omterm{def:g8:puberty:puberty}{puberty}, and name the stage of
\cref{def:g3:stages-of-life:stages} whose core it is.
\end{exercise}

\begin{exercise}[$\star$]\label{exo:g8:puberty:2}
List three changes of girls' \omterm{def:g8:puberty:puberty}{puberty}, three of boys',
and two common to both.
\end{exercise}

\begin{exercise}[$\star$]\label{exo:g8:puberty:3}
What marks the start of \omterm{def:g8:sexual-reproduction:gametes}{gamete} production in each sex,
by \cref{prop:g8:puberty:changes}?
\end{exercise}

\begin{exercise}[$\star$]\label{exo:g8:puberty:4}
What is a \omterm{def:g8:puberty:hormones}{hormone} --- released where, carried how,
obeyed by whom?
\end{exercise}

\begin{exercise}[$\star$]\label{exo:g8:puberty:5}
Where does \omterm{def:g8:puberty:puberty}{puberty}'s chain of command begin, and where
does it end?
\end{exercise}

\begin{exercise}[$\star$]\label{exo:g8:puberty:6}
Two thirteen-year-olds stand at visibly different
stages. What does \cref{rem:g8:puberty:timetable} say
--- in full?
\end{exercise}

\begin{exercise}[$\star\star$]\label{exo:g8:puberty:7}
Why can one messenger coordinate building sites half a
body apart, where \omterm{def:g5:brain-in-command:system}{nerves} would need wiring to each? Use
\cref{ex:g8:puberty:sites}.
\end{exercise}

\begin{exercise}[$\star\star$]\label{exo:g8:puberty:8}
Explain the mood weather of \cref{rem:g8:puberty:moods}
--- two causes, one reassurance.
\end{exercise}

\begin{exercise}[$\star\star$]\label{exo:g8:puberty:9}
Connect the growth burst to the plate: which \omterm{def:g7:digestion-and-nutrients:families}{nutrient}
families does the rebuild draw hardest, and from which
chapter's reasoning?
\end{exercise}

\begin{exercise}[$\star\star$]\label{exo:g8:puberty:10}
Why do \cref{prop:g5:healthy-choices:substances}'s
substances strike an \omterm{def:g3:stages-of-life:stages}{adolescent} body hardest? (The
construction image carries the answer.)
\end{exercise}

\begin{exercise}[$\star\star$]\label{exo:g8:puberty:11}
Sort into trustworthy and not, with one reason each:
the school nurse; a viral video's miracle claims; this
book; playground legend; the family doctor.
\end{exercise}

\begin{exercise}[$\star\star\star$]\label{exo:g8:puberty:12}
``\omterm{def:g8:puberty:puberty}{Puberty} is the \omterm{def:g8:puberty:hormones}{hormone} system's public debut.''
Defend the sentence: what does the rebuild demonstrate
about that system's reach, coordination and difference
from \omterm{def:g5:brain-in-command:system}{nerve} signaling --- three claims, three
observations.
\end{exercise}

\section{Problem: The Growth-Chart Conference}

\begin{problem}[{Weekend problem --- four (invented)
growth charts, read with the whole chapter}]\label{pb:g8:puberty:1}
A school doctor reviews four charts, height against age,
each with its measuring dates: L, whose curve leapt
steeply from eleven to thirteen and is now flattening at
fifteen; M, flat-climbing at thirteen, burst not yet
begun; N, mid-burst at fourteen; O, whose steep segment
ran nine to twelve and has plateaued.

\medskip\noindent\textbf{Part I --- Reading the curves.}
\begin{enumerate}
  \item Place each chart on
        \cref{rem:g3:stages-of-life:bursts}'s map: which
        burst is showing, and where does each student
        stand on it?
  \item M worries about being ``left behind''. Draft the
        doctor's answer from
        \cref{rem:g8:puberty:timetable}, with the
        prediction the chart's flat-climb supports.
  \item O asks whether the plateau means ``done growing
        at twelve''. Correct the reading: what does the
        plateau end, and what continues after it?
  \item Why does the doctor read the \emph{slope}, not
        the height, to judge where each stands?
\end{enumerate}

\medskip\noindent\textbf{Part II --- Behind the curves.}
\begin{enumerate}[resume]
  \item Name the machinery that started each burst, and
        the route its orders traveled.
  \item The doctor also tracks \omterm{prop:g1:healthy-habits:needs}{sleep} and diet at these
        visits. Tie each to its line of
        \cref{prop:g8:puberty:care}.
  \item N reports mood storms and a late-drifting
        bedtime. Give the two-part explanation of
        \cref{rem:g8:puberty:moods} --- and the practical
        consequence for N's school-week \omterm{prop:g1:healthy-habits:needs}{sleep}.
  \item L asks the doctor privately about a change the
        chart cannot show. State, without inventing the
        question, why the consulting room is the right
        address (\cref{ex:g8:puberty:questions}).
\end{enumerate}

\medskip\noindent\textbf{Part III --- The conference
note.}
\begin{enumerate}[resume]
  \item Write the one-line note for each chart: stage,
        normality, and any care point.
  \item The four charts differ by years, yet the doctor
        files all four under ``normal''. Justify in one
        sentence.
  \item A parent asks for ``something to speed M up''.
        Answer with the timetable remark and the
        machinery: what is there to fix?
  \item Close with the chapter's one-sentence summary
        for the class notice board: what \omterm{def:g8:puberty:puberty}{puberty} is,
        who runs it, and what it needs.
\end{enumerate}
\end{problem}
